% Corpo do artigo no template de TCC de SI Contagem (PUC Minas).
% O título, os autores, o resumo e o preâmbulo ficam no arquivo principal do template.
% Título: A Implementação de um Mapa Interativo Sobre as Empresas de TI do Município de Contagem
% As citações dependem das entradas de referencias.bib e do arquivo abntex2-alf.bst no projeto.
% A Figura 1 depende de \usepackage{tikz} e \usetikzlibrary{arrows.meta,positioning} no preâmbulo.

\section{\esp Introdução}

O uso de dados para compreender a dinâmica econômica de um município pode contribuir para a identificação de características relacionadas à distribuição das atividades econômicas, à ocupação territorial e ao desenvolvimento de determinados setores.
Nesse contexto, dados produzidos e mantidos por órgãos públicos podem ser utilizados não apenas para fins administrativos, mas também para produzir informações que auxiliem cidadãos, empresas, pesquisadores e gestores na compreensão do ambiente econômico local.
Parte desses registros é publicada como dado aberto, em arquivos que cobrem o Brasil ou uma unidade da federação e que podem ser filtrados até o município.
Nenhuma dessas fontes, isoladamente, reúne a identificação do estabelecimento, a atividade econômica e a coordenada geográfica.
Assim, conhecer um setor em uma cidade exige cruzar dados de mais de uma instituição.

A existência de dados abertos com essas características oferece diferentes possibilidades de análise. Porém, o arquivo original, do jeito que o órgão publica, não é simples de consultar.
Neste trabalho, acessível quer dizer facilidade de uso: os dados ficam reunidos em um único site, com uma interface clara, em que a pessoa vê mapa, gráfico e tabela sem precisar baixar várias bases e cruzá-las.
Não se trata de acessibilidade digital no sentido de normas para pessoas com deficiência, como a WCAG ou o Modelo de Acessibilidade em Governo Eletrônico (eMAG).
Os arquivos podem conter muitos registros, códigos e campos de endereço, espalhados em mais de uma fonte, o que dificulta a leitura quando eles são abertos direto na tabela.
Segundo \citeonline{ansari2022}, mesmo com mais dados governamentais disponíveis, o uso ainda pode ser limitado, e a visualização pode deixar essas informações mais fáceis de entender.
Diante desse cenário, este trabalho busca responder à seguinte questão de pesquisa: como dados abertos podem ser combinados para caracterizar e apresentar, de maneira acessível e interativa, informações sobre as empresas do setor de Tecnologia da Informação localizadas no município de Contagem?
Além de descrever essas empresas, o trabalho observa onde elas estão no município, a partir do endereço cruzado com bases geográficas abertas.
A análise não fica só no mapa: também entra o perfil dos estabelecimentos, como atividade e porte.

Neste trabalho, empresa de TI é o estabelecimento com situação cadastral ativa em Contagem cuja atividade principal está na divisão 62 da Classificação Nacional de Atividades Econômicas (CNAE), a de serviços de tecnologia da informação.
Entram desenvolvimento de software, web design, consultoria e suporte em TI.
Ficam de fora telecomunicações, hospedagem de dados, comércio e fabricação de equipamentos, conserto de computador e curso de informática.
A lista fechada de códigos está na metodologia.
A unidade observada é o estabelecimento do CNPJ, e não o grupo econômico: matriz e filial contam separado quando as duas estão no município e na atividade escolhida.

Para responder a essa questão, o objetivo geral deste trabalho é desenvolver uma aplicação web interativa que permita caracterizar e visualizar as empresas do setor de Tecnologia da Informação do município de Contagem a partir da integração de dados abertos.
A aplicação reunirá indicadores, gráficos, mapas, tabelas e mecanismos de filtragem, ajustados à estrutura e à qualidade das fontes utilizadas.

Como objetivos específicos, pretende-se:

\begin{itemize}
    \item coletar e tratar dados abertos sobre os estabelecimentos de Tecnologia da Informação de Contagem, integrando atividade econômica, porte e localização em uma base única;
    \item projetar uma aplicação web para a consulta desses dados, com levantamento de requisitos, diagramas de casos de uso e de classes e protótipo das interfaces;
    \item implementar essa aplicação, de modo que o usuário possa consultar, filtrar e interpretar as informações por meio de mapa, gráficos e tabelas;
    \item testar a aplicação, conferindo a localização atribuída aos estabelecimentos e avaliando a usabilidade da interface por meio de inspeção heurística.
\end{itemize}

Os dados usados são cadastrais e já publicados. Ainda assim, o trabalho leva em conta a Lei Geral de Proteção de Dados Pessoais, Lei nº 13.709, de 14 de agosto de 2018 \cite{brasil2018}.
A aplicação não carrega a base de sócios, em que aparece CPF de pessoa física, e não pede cadastro de quem consulta o site.
O que se mostra é dado do estabelecimento. O detalhe desse cuidado está na metodologia.

A relevância deste estudo justifica-se pela necessidade de aproximar dados públicos, hoje dispersos e difíceis de consultar na forma bruta, de públicos que podem utilizá-los na prática.
Os públicos previstos são:

\begin{itemize}
    \item cidadãos poderão conhecer melhor as empresas de tecnologia existentes no município;
    \item profissionais poderão identificar organizações relacionadas às suas áreas de interesse;
    \item empresários poderão utilizar as informações para conhecer empresas do mesmo setor e identificar possíveis fornecedores ou parceiros;
    \item empreendedores poderão explorar características da distribuição das atividades econômicas no município.
\end{itemize}

Sob a perspectiva acadêmica, o trabalho contribui para a área de Sistemas de Informação ao integrar dados abertos, análise descritiva e visualização interativa em um recorte municipal concreto, alinhado a estudos que relacionam localização e setor de tecnologia em contextos urbanos, como o de \citeonline{zenka2021}.
A escolha do tema pauta-se na formação em Sistemas de Informação e no interesse por soluções que transformem bases públicas em instrumentos de consulta e interpretação para usuários não especializados.

Este trabalho está estruturado em seis seções.
Após esta introdução, a Seção 2 apresenta o referencial teórico, e a Seção 3 os trabalhos relacionados.
A Seção 4 descreve a metodologia empregada, incluindo a coleta e a integração das bases de dados abertos e o projeto e a implementação da aplicação web.
A Seção 5 apresenta e discute os resultados obtidos, bem como os testes da aplicação.
Por fim, a Seção 6 apresenta as conclusões do trabalho, suas limitações e as possibilidades de continuidade da pesquisa.

\section{\esp Referencial Teórico}

Esta seção apresenta os conceitos que fundamentam o trabalho.
São tratados, em sequência, dados abertos, visualização e usabilidade, avaliação heurística, a classificação das atividades econômicas, a distribuição geográfica das empresas e as aplicações web de consulta.

\subsection{\esp Dados abertos governamentais}

Dados abertos governamentais são informações produzidas ou mantidas pelo poder público e disponibilizadas para que outras pessoas possam usá-las de novo.
O objetivo costuma ser aumentar a transparência e permitir que esses dados sejam reaproveitados.
\citeonline{molodtsov2024} estudaram a usabilidade de portais de dados abertos em 33 países.
O quadro que eles aplicam não olha só se o arquivo está disponível para \textit{download}.
Também entra se a pessoa consegue explorar o dado e entendê-lo, mesmo sem prática com planilha.
\citeonline{ansari2022}, ao revisar pesquisas sobre visualização de dados governamentais, observam um ponto parecido: publicar o arquivo não faz, por si só, com que o dado seja usado.
Muita gente não consegue ler o formato, cruzar bases diferentes e interpretar a tabela.
Uma aplicação que junta as bases e mostra mapa, gráfico e filtro pode tornar essa consulta mais simples.
Esse é o caminho deste trabalho: caracterizar as empresas de TI de Contagem a partir de bases abertas nacionais, e não a partir de um arquivo municipal já pronto.

\subsection{\esp Visualização de dados e usabilidade}

A visualização de dados é uma forma de reduzir a distância entre o arquivo publicado e o uso por quem não tem domínio técnico.
\citeonline{mesquita2026} analisaram 30 painéis e 53 telas de consulta de dados de orçamento público no Brasil.
O achado principal foi que a maior parte dos portais entrega o dado, mas pouco ajuda quem não tem familiaridade com números e tabelas.
Falta contexto, e a tela quase não considera a pessoa que está consultando.
\citeonline{ansari2022} descrevem uma lacuna parecida: há muitos relatos de visualização de dados abertos e pouca avaliação de usabilidade para públicos diferentes.
Para os autores, gráfico, mapa e consulta junto do dado tendem a ser mais úteis do que a entrega isolada do arquivo.
Por isso este trabalho reúne indicadores, tabelas, filtros e mapa, em vez de disponibilizar só a planilha tratada.

\subsection{\esp Avaliação heurística de usabilidade}

Usabilidade, neste trabalho, é a facilidade de usar a interface.
A pessoa encontra a informação, entende o que está vendo e consegue filtrar sem se perder.
Como já foi dito na introdução, não se trata de normas de acessibilidade digital.
A avaliação heurística é uma inspeção da tela.
Quem avalia percorre a interface com uma lista de princípios e anota os problemas encontrados.
A lista mais usada é a das dez heurísticas de \citeonline{nielsen1994}.
Ela inclui, entre outros pontos, deixar claro o que o sistema está fazendo, manter um padrão visual, evitar erro e oferecer ajuda quando algo falha.
\citeonline{vazquezingelmo2024} aplicaram essas heurísticas em uma ferramenta de criação de visualizações.
Cada avaliador registrou o problema, a heurística relacionada e a gravidade.
O estudo mostra que o método ajuda a localizar falhas de navegação, de retorno ao usuário e de organização da tela antes de um teste com o público final.
Este trabalho usa o mesmo tipo de inspeção na aplicação desenvolvida, para registrar o que dificulta a consulta ao mapa, aos filtros, aos gráficos e à tabela.

\subsection{\esp Classificação de atividades econômicas e integração de bases}

A identificação de um setor em cadastros públicos depende de uma classificação padronizada das atividades.
No Brasil, a Classificação Nacional de Atividades Econômicas (CNAE), mantida pelo IBGE, é o instrumento usado pelo sistema estatístico e pelos órgãos que registram empresas \cite{ibge2020}.
A CNAE organiza as atividades em níveis, da seção até a subclasse.
No CNPJ, cada estabelecimento tem um código de atividade principal.

Neste trabalho, empresa de TI é o estabelecimento cuja atividade principal está na divisão 62, de serviços de tecnologia da informação, na versão CNAE-Subclasses 2.3.
Os códigos são 6201-5/01 (desenvolvimento de programas sob encomenda), 6201-5/02 (web design), 6202-3/00 (programas customizáveis), 6203-1/00 (programas não customizáveis), 6204-0/00 (consultoria em tecnologia da informação) e 6209-1/00 (suporte técnico e outros serviços em TI).
Não entram a divisão 63 (hospedagem, tratamento de dados e portais), o comércio e a fabricação de equipamentos, o conserto de computador nem o curso de informática.
A regra usa só a atividade principal, para a lista não crescer com códigos secundários de empresas de outro ramo.
O código, sozinho, não monta o mapa.
É preciso juntar o cadastro de estabelecimentos, a tabela da CNAE, o endereço com coordenada e o contorno do município.
Sem essa junção, a atividade, o endereço e a localização ficam em arquivos separados.

\subsection{\esp Distribuição espacial de empresas e geovisualização}

A localização da empresa no município é parte da análise, e não só um campo de cadastro.
\citeonline{zenka2021}, ao estudar as firmas de tecnologia da informação e comunicação em Ostrava, na Tchéquia, encontraram concentração no centro e em áreas mais internas da cidade, e não uma dispersão uniforme.
Contagem não é Ostrava, e o método de lá foi mapeamento em campo.
Ainda assim, a contribuição que interessa aqui permanece: olhar a posição dentro da cidade muda a leitura do setor, coisa que um total municipal não mostra.
\citeonline{ansari2022} incluem o mapa entre os recursos que ajudam a entender conjuntos de dados governamentais.
A coordenada deste trabalho é estimada a partir do endereço e de bases territoriais abertas.
Ela não substitui uma medição no terreno.
Ela alimenta o mapa de pontos e a contagem por área do município.
Como essas áreas não têm o mesmo tamanho, a comparação entre elas também usa uma medida relativa, como o número de estabelecimentos por quilômetro quadrado, e não apenas o total bruto.

\subsection{\esp Aplicações web interativas para exploração de dados}

Na prática, esses conceitos aparecem em uma aplicação que junta consulta, filtro e visualização.
\citeonline{ansari2022} recomendam plataformas em que o gráfico e a consulta já venham com o dado, e em que o uso seja avaliado.
\citeonline{mesquita2026} mostram que, no Brasil, muitos painéis públicos de orçamento ainda não chegam nesse ponto: o acesso existe, mas a tela não ajuda quem tem pouca prática com dados.
Essas referências apoiam dois objetivos deste trabalho: implementar a aplicação e testar a usabilidade dela por inspeção heurística.
Na aplicação, indicadores, gráficos, mapa e tabelas funcionam sobre a mesma base, para a pessoa explorar as empresas de TI de Contagem com os atributos que as fontes permitem.

\section{\esp Trabalhos Relacionados}

\section{\esp Metodologia}

Esta pesquisa tem como objetivo desenvolver uma aplicação web interativa que caracterize e apresente as empresas de Tecnologia da Informação do município de Contagem a partir da integração de dados abertos.
Para isso, bases públicas sobre os estabelecimentos, suas atividades econômicas e seus endereços são coletadas, tratadas e combinadas em um banco de dados único.
Esse banco alimenta uma aplicação com mapa, gráficos, tabelas e filtros, cuja usabilidade é avaliada por inspeção heurística.

\subsection{\esp Classificação da pesquisa}

Quanto aos objetivos, a pesquisa classifica-se como descritiva.
Segundo \citeonline{gil2008}, a pesquisa descritiva registra características de determinada população ou fenômeno e permite a organização dessas informações em tabelas, gráficos e outras formas de apresentação.
No presente trabalho, descrevem-se as empresas de Tecnologia da Informação localizadas em Contagem, a partir de atributos cadastrais e da distribuição geográfica observada nas bases abertas.

Quanto aos procedimentos, trata-se de pesquisa documental.
Utilizam-se documentos e bases oficiais já existentes, publicadas por órgãos públicos, e não a coleta direta de entrevistas ou questionários com pessoas.

Quanto à abordagem, adota-se o caráter misto.
A dimensão quantitativa aparece na contagem de estabelecimentos, na distribuição por subclasses da CNAE, por porte e por recortes territoriais.
A dimensão qualitativa aparece na interpretação dos padrões espaciais, na inspeção heurística da interface e na discussão das limitações das fontes, sem pretender estabelecer relações causais entre localização e desempenho econômico.

\subsection{\esp Objeto de estudo e recorte}

O objeto de estudo são os estabelecimentos de Tecnologia da Informação situados no município de Contagem (MG), identificados pelo código IBGE 3118601.
O setor fica delimitado pela CNAE-Subclasses 2.3, só com a atividade principal na divisão 62.
Os códigos incluídos são:

\begin{itemize}
    \item 6201-5/01, desenvolvimento de programas de computador sob encomenda;
    \item 6201-5/02, web design;
    \item 6202-3/00, desenvolvimento e licenciamento de programas de computador customizáveis;
    \item 6203-1/00, desenvolvimento e licenciamento de programas de computador não customizáveis;
    \item 6204-0/00, consultoria em tecnologia da informação;
    \item 6209-1/00, suporte técnico, manutenção e outros serviços em tecnologia da informação.
\end{itemize}

Ficam de fora, neste recorte, hospedagem e portais (divisão 63), comércio e fabricação de equipamentos de informática, conserto de computador e curso de informática.
A unidade de análise é o estabelecimento do cadastro de CNPJ.
Matriz e filial entram como registros distintos quando estiverem no município e na atividade selecionada.
Serão considerados estabelecimentos com situação cadastral ativa na fotografia da base utilizada.
O CNAE secundário não será usado como critério de entrada.

A descrição dos dados será orientada por quatro perguntas: onde esses estabelecimentos se concentram em Contagem; qual é o perfil de porte; qual é a proporção de microempreendedores individuais; e qual parte dos endereços recebe coordenada em cada nível de junção.

Não haverá teste com usuários nem questionário.
A avaliação da interface será uma inspeção feita pelo próprio autor, descrita na Seção \ref{subsec:testes}.
Por isso, não se aplica tramitação em comitê de ética para coleta com pessoas.
Os dados empregados são públicos e abertos.

\subsection{\esp Proteção de dados}

A Lei nº 13.709/2018, a Lei Geral de Proteção de Dados Pessoais, trata de dado de pessoa natural \cite{brasil2018}.
O foco deste trabalho é o estabelecimento, isto é, a unidade cadastrada no CNPJ.
A base de sócios da Receita Federal, que pode trazer CPF de pessoa física, mesmo mascarado, não será baixada nem carregada na aplicação.
Telefone e e-mail do estabelecimento também não entram na tela, porque não são necessários para descrever o setor nem para montar o mapa.
O que será exibido é informação que a Receita já publica como dado aberto do estabelecimento: nome, CNAE, porte, situação cadastral e endereço.
No caso do microempreendedor individual, o nome do cadastro pode ser o nome da pessoa.
Esse nome já é público no CNPJ.
A aplicação não junta esse nome com CPF nem com quadro de sócios.
Não haverá cadastro de usuário no site.

\subsection{\esp Etapas do trabalho}

Esta pesquisa será realizada em quatro etapas: coleta e tratamento dos dados, desenvolvimento do projeto da solução, implementação da solução e testes.
Na primeira etapa, os estabelecimentos de TI de Contagem são extraídos das bases abertas e integrados em um banco de dados local, com atividade, porte e coordenada.
Na segunda, a aplicação é projetada por meio do levantamento de requisitos, dos diagramas de casos de uso e de classes e do protótipo das interfaces.
Na terceira, a aplicação é implementada sobre esse banco.
Na quarta, são conferidas a localização atribuída aos estabelecimentos e a usabilidade da interface.
A Figura \ref{fig:fluxo-metodologia} resume esse fluxo.

% Fluxograma da metodologia — incluído em trabalho_final.tex por % Fluxograma da metodologia — incluído em trabalho_final.tex por % Fluxograma da metodologia — incluído em trabalho_final.tex por \input{figuras/fluxo-metodologia}
% Requer no preâmbulo: \usepackage{tikz} e \usetikzlibrary{arrows.meta,positioning}
\begin{figure}[htbp]
    \centering
    \caption{Fluxo geral das etapas metodológicas}
    \label{fig:fluxo-metodologia}
    \vspace{-0.2cm}
    \begin{tikzpicture}[
        node distance=0.55cm,
        etapa/.style={
            draw=black,
            thick,
            rounded corners=2pt,
            align=center,
            text width=9.2cm,
            inner sep=6pt,
            font=\small,
            fill=white
        },
        seta/.style={-Latex, thick}
    ]
        \node[etapa] (a) {\textbf{1. Coleta e tratamento dos dados}\\ \footnotesize CNPJ, CNAE, CNEFE e malhas do IBGE integrados em banco local};
        \node[etapa, below=of a] (b) {\textbf{2. Projeto da solução}\\ \footnotesize requisitos, casos de uso, classes e protótipo das interfaces};
        \node[etapa, below=of b] (c) {\textbf{3. Implementação da solução}\\ \footnotesize programas de carga e aplicação web};
        \node[etapa, below=of c] (d) {\textbf{4. Testes}\\ \footnotesize conferência da localização, comparação com o CEMPRE e inspeção heurística};

        \draw[seta] (a) -- (b);
        \draw[seta] (b) -- (c);
        \draw[seta] (c) -- (d);
    \end{tikzpicture}
    \vspace{-0.2cm}
    \\\textbf{\footnotesize Fonte: Elaborado pelo autor.}
\end{figure}

% Requer no preâmbulo: \usepackage{tikz} e \usetikzlibrary{arrows.meta,positioning}
\begin{figure}[htbp]
    \centering
    \caption{Fluxo geral das etapas metodológicas}
    \label{fig:fluxo-metodologia}
    \vspace{-0.2cm}
    \begin{tikzpicture}[
        node distance=0.55cm,
        etapa/.style={
            draw=black,
            thick,
            rounded corners=2pt,
            align=center,
            text width=9.2cm,
            inner sep=6pt,
            font=\small,
            fill=white
        },
        seta/.style={-Latex, thick}
    ]
        \node[etapa] (a) {\textbf{1. Coleta e tratamento dos dados}\\ \footnotesize CNPJ, CNAE, CNEFE e malhas do IBGE integrados em banco local};
        \node[etapa, below=of a] (b) {\textbf{2. Projeto da solução}\\ \footnotesize requisitos, casos de uso, classes e protótipo das interfaces};
        \node[etapa, below=of b] (c) {\textbf{3. Implementação da solução}\\ \footnotesize programas de carga e aplicação web};
        \node[etapa, below=of c] (d) {\textbf{4. Testes}\\ \footnotesize conferência da localização, comparação com o CEMPRE e inspeção heurística};

        \draw[seta] (a) -- (b);
        \draw[seta] (b) -- (c);
        \draw[seta] (c) -- (d);
    \end{tikzpicture}
    \vspace{-0.2cm}
    \\\textbf{\footnotesize Fonte: Elaborado pelo autor.}
\end{figure}

% Requer no preâmbulo: \usepackage{tikz} e \usetikzlibrary{arrows.meta,positioning}
\begin{figure}[htbp]
    \centering
    \caption{Fluxo geral das etapas metodológicas}
    \label{fig:fluxo-metodologia}
    \vspace{-0.2cm}
    \begin{tikzpicture}[
        node distance=0.55cm,
        etapa/.style={
            draw=black,
            thick,
            rounded corners=2pt,
            align=center,
            text width=9.2cm,
            inner sep=6pt,
            font=\small,
            fill=white
        },
        seta/.style={-Latex, thick}
    ]
        \node[etapa] (a) {\textbf{1. Coleta e tratamento dos dados}\\ \footnotesize CNPJ, CNAE, CNEFE e malhas do IBGE integrados em banco local};
        \node[etapa, below=of a] (b) {\textbf{2. Projeto da solução}\\ \footnotesize requisitos, casos de uso, classes e protótipo das interfaces};
        \node[etapa, below=of b] (c) {\textbf{3. Implementação da solução}\\ \footnotesize programas de carga e aplicação web};
        \node[etapa, below=of c] (d) {\textbf{4. Testes}\\ \footnotesize conferência da localização, comparação com o CEMPRE e inspeção heurística};

        \draw[seta] (a) -- (b);
        \draw[seta] (b) -- (c);
        \draw[seta] (c) -- (d);
    \end{tikzpicture}
    \vspace{-0.2cm}
    \\\textbf{\footnotesize Fonte: Elaborado pelo autor.}
\end{figure}


\subsection{\esp Coleta e tratamento dos dados}
\label{subsec:coleta}

A coleta será feita em bases abertas oficiais, por meio do \textit{download} e da consulta aos arquivos publicados por órgãos públicos.

A principal fonte da lista de estabelecimentos será a base pública de CNPJ da Receita Federal do Brasil, que reúne identificação, CNAE, situação cadastral e endereço declarado.
O arquivo de sócios dessa base não será utilizado.

A estrutura e a descrição das atividades virão da CNAE-Subclasses 2.3, mantida pelo IBGE, usada para fechar o recorte da divisão 62.

As coordenadas geográficas serão obtidas no Cadastro Nacional de Endereços para Fins Estatísticos (CNEFE), do IBGE, associado ao Censo Demográfico de 2022.
A base inteira não será geocodificada em lote por serviço online: o CNEFE é a fonte da coordenada.
As malhas territoriais do IBGE fornecerão o contorno do município e, quando disponíveis no nível adequado, recortes internos úteis à agregação espacial, com a área em quilômetros quadrados para a medida relativa.

Para a conferência dos totais, serão consultadas tabelas do Cadastro Central de Empresas (CEMPRE) por meio do SIDRA, sempre no município de Contagem e em nível de CNAE compatível com o recorte.

Como apoio à extração, poderá ser utilizada a plataforma Base dos Dados, desde que os filtros reproduzam o mesmo universo da fonte oficial.

Os arquivos nacionais não serão carregados integralmente na aplicação.
O tratamento ocorrerá em ambiente local, com filtragem prévia para Contagem e para as subclasses da divisão 62.

A integração seguirá chaves explícitas: código de município da Receita para isolar Contagem; CNPJ básico para unir estabelecimentos e dados da empresa, como razão social, porte e capital social; CNAE sem máscara para padronizar a atividade; código IBGE 3118601 para ligar o resultado às malhas, ao CNEFE e ao CEMPRE.

O endereço passa por uma padronização antes da junção com o CNEFE.
O CEP fica com 8 dígitos, sem hífen.
O texto vai para maiúsculas, sem acento.
Abreviações comuns de tipo de logradouro são expandidas, como R.\ para RUA e AV.\ para AVENIDA, nos dois lados da junção.
Número vazio ou registrado como S/N não entra na tentativa mais específica.

A junção com o CNEFE segue uma ordem.
Primeiro, tenta-se CEP, logradouro padronizado e número.
Se não houver correspondência, tenta-se só o CEP, e o ponto é marcado como aproximado.
Se o CEP do CNPJ não existir no CNEFE, o estabelecimento continua na base para as contagens e as tabelas, mas fica de fora do mapa de pontos.
Cada registro guarda o nível obtido: endereço completo, só CEP ou sem coordenada.
Com isso, calcula-se a taxa de acerto de cada nível.

O registro também guarda o nível de geocodificação informado pelo próprio CNEFE, na variável \texttt{NV\_GEO\_COORD}.
Esse campo indica se a coordenada é a original do Censo 2022, se foi modificada para apartamentos no mesmo número, se foi estimada pelo IBGE ou se foi atribuída por face de quadra, localidade ou setor censitário.
Assim, uma junção por endereço completo que herdou uma coordenada estimada pelo IBGE não é contada como ponto coletado em campo.

Quando várias empresas recebem a mesma coordenada, o que é esperado no nível de CEP, essa sobreposição fica registrada.
Assim, um CEP com muitas empresas não é tratado como um único endereço conferido no terreno.

O resultado dessa integração será gravado em um banco de dados local, com a data da fotografia das bases.
A aplicação consulta esse banco, e não os arquivos nacionais nem os sites da Receita e do IBGE.

\subsection{\esp Aplicação desenvolvida}

A aplicação será documentada segundo as etapas usuais de projeto de software: análise de requisitos, diagramas de casos de uso e de classes e protótipo das interfaces.
A arquitetura e as ferramentas da implementação também são descritas nesta seção.

\subsubsection{\esp Projeto da solução}

Os requisitos funcionais previstos são:

\begin{itemize}
    \item RF01: exibir indicadores gerais, como o total de estabelecimentos ativos, a distribuição por subclasse da CNAE, o porte e a proporção de microempreendedores individuais;
    \item RF02: filtrar os estabelecimentos por subclasse da CNAE, porte, condição de microempreendedor individual e área do município;
    \item RF03: exibir os estabelecimentos em um mapa de pontos, indicando quantos estão em uma mesma coordenada;
    \item RF04: exibir a quantidade de estabelecimentos por área do município, em total e por quilômetro quadrado;
    \item RF05: listar os estabelecimentos em uma tabela com busca por nome, com CNAE, porte, situação cadastral e endereço;
    \item RF06: informar a data da fotografia das bases e o nível de coordenada de cada estabelecimento.
\end{itemize}

Os requisitos não funcionais previstos são:

\begin{itemize}
    \item RNF01: acesso pelo navegador, sem cadastro de usuário;
    \item RNF02: não exibir CPF, quadro de sócios, telefone ou e-mail;
    \item RNF03: indicadores, gráficos, mapa e tabela devem responder aos mesmos filtros, sobre os mesmos registros;
    \item RNF04: as consultas devem ser feitas no banco local, sem depender de serviço externo além do mapa de fundo.
\end{itemize}

O diagrama de casos de uso terá dois atores.
O visitante consulta os indicadores, aplica filtros, explora o mapa e consulta a tabela.
O administrador, papel exercido pelo autor, carrega uma nova fotografia das bases.

O diagrama de classes representará o estabelecimento, a atividade econômica, a localização, com a coordenada e o nível de junção, a área do município, com a área em quilômetros quadrados, e a carga de dados, com a data da fotografia e a versão das bases.

O protótipo das interfaces cobrirá o painel de indicadores com os filtros, o mapa e a tabela.
Os diagramas e o protótipo serão apresentados como figuras nesta seção.

\subsubsection{\esp Implementação da solução}

A arquitetura terá três partes: os programas de carga, o banco de dados local e a aplicação web.
Os programas de carga leem as bases, aplicam os filtros e as junções descritos na Seção \ref{subsec:coleta} e gravam o resultado no banco.
A aplicação web apenas consulta esse banco.
Uma atualização dos dados será uma nova carga, com outra data de fotografia, e não uma consulta ao vivo na hora em que alguém abre o site.

Para o tratamento dos dados, serão utilizadas ferramentas de análise e manipulação tabular em linguagem Python, com apoio de bibliotecas adequadas à leitura de arquivos grandes e à filtragem por município e CNAE.

Para a interface, serão empregados recursos de desenvolvimento web e uma biblioteca de mapas interativos.
O mapa de fundo utilizará \textit{tiles} do OpenStreetMap.

As versões das bases, as datas de \textit{download} e as regras de filtro serão registradas em cada carga, para que o processo possa ser refeito.

\subsection{\esp Testes}
\label{subsec:testes}

Os testes verificam dois aspectos da aplicação: a localização atribuída aos estabelecimentos e a usabilidade da interface.

Para estimar o erro da localização, uma amostra de cerca de 30 endereços será conferida em outra fonte, por consulta manual no mapa ou por um serviço de geocodificação usado só nessa amostra.
A conferência não substitui o CNEFE na base inteira.
Ela serve para dizer, de forma aproximada, o quanto o ponto se afasta do endereço informado.

Os totais da base também serão comparados com os do CEMPRE.
A comparação não busca números iguais, porque o ano e os critérios de cada base são diferentes.
O que se verifica é se a ordem de grandeza é coerente.

A usabilidade será avaliada por inspeção heurística, feita pelo autor, sem usuários convidados.
A lista usada é a das dez heurísticas de \citeonline{nielsen1994}, no mesmo tipo de registro adotado por \citeonline{vazquezingelmo2024}: problema encontrado, heurística relacionada e gravidade.
A gravidade será anotada em três níveis: leve, médio e grave.
As telas percorridas serão o mapa, os filtros, os gráficos e a tabela.
O resultado dessa inspeção entra na seção de resultados, como lista de problemas e ajustes possíveis, e não como nota da aplicação.

\subsection{\esp Procedimentos de análise}

A análise será descritiva e seguirá as quatro perguntas do objeto de estudo.
Serão calculados totais de estabelecimentos ativos no recorte, a distribuição por subclasse da CNAE, o perfil de porte e a proporção de microempreendedores individuais.
Também serão apresentados os percentuais de georreferenciamento em cada nível de junção.

No mapa, serão observados onde os pontos se concentram, onde há vazios e como os estabelecimentos se distribuem pelos recortes territoriais disponíveis.
A comparação entre áreas usa o total e uma medida normalizada, o número de estabelecimentos por quilômetro quadrado, para uma área maior não parecer mais concentrada só por ter mais espaço.
Mapas de pontos e, quando a taxa de coordenada permitir, uma representação de densidade ajudam essa leitura, com o aviso sobre pontos sobrepostos.

Os resultados serão apresentados em tabelas, gráficos e mapas, e discutidos na seção de resultados.

\subsection{\esp Limitações metodológicas previstas}

Algumas limitações vêm das fontes e serão descritas com os resultados.
O CNAE é a atividade declarada no cadastro, não uma verificação do que a empresa faz no dia a dia.
Usar só a atividade principal deixa de fora estabelecimentos de outro ramo que também prestam serviço de TI.

A coordenada do CNEFE foi registrada em campo pelo IBGE durante o Censo de 2022, e o próprio cadastro informa o nível de precisão de cada ponto.
Ela representa o endereço censitário, não a localização auditada do estabelecimento.

A principal fonte de erro está na junção.
Como o CNEFE não tem CNPJ, a ligação é feita pelo endereço.
Quando apenas o CEP coincide, a coordenada atribuída é uma aproximação, e vários estabelecimentos podem cair no mesmo ponto.

A padronização do endereço, a ordem da junção, a taxa de acerto de cada nível e a amostra de cerca de 30 endereços existem para deixar esse erro visível.

Porte e capital social estão no CNPJ.
O número de pessoas ocupadas por empresa, em geral, não está disponível de forma pública e identificável nas bases usadas aqui.
Totais do CEMPRE não substituem a lista de CNPJs no mapa e não devem ser lidos como a mesma contagem.
O banco guarda a fotografia do mês carregado. O CNPJ muda todo mês, e a aplicação não acompanha essa mudança sozinha.
A inspeção heurística, feita por uma pessoa só, encontra problemas de uso, mas não mede como um grupo de usuários de fato utiliza a tela.
Esses limites indicam o que os resultados podem afirmar.

\section{\esp Resultados}

\section{\esp Considerações Finais}
